\chapter{The A* Search Framework for MLP Training}
\label{chap:astar_framework}

The primary objective of this project is to reformulate the optimization challenge of training a Multilayer Perceptron (MLP) as a \textbf{shortest-path problem} within a discrete, high-dimensional weight space. This chapter introduces the framework designed to achieve this, detailing how the continuous weight space is discretized, how the nodes and edges of the search graph are defined, and how the A* algorithm's evaluation function is adapted to guide the network towards optimal weights.

\section{Discretization of the Weight Space}

The A* algorithm fundamentally operates on a discrete graph. Therefore, the continuous domain of real-valued weight parameters must be converted into a finite, discrete set of possible states. This is handled by the \texttt{QuantizedMLP} class.

\subsection{Quantized State Definition}
To define the search nodes, the floating-point parameters (weights and biases) of the MLP are subjected to uniform quantization. Given a predefined \textbf{quantization factor} ($Q$), every weight $w_i$ is rounded to the nearest multiple of $1/Q$.

$$w_{i, \text{quantized}} = \text{round}(w_i \cdot Q) / Q$$

The number of potential discrete weight configurations defines the size of the search space. To manage the complexity and avoid instability, the parameters are constrained to a defined numerical range $[\text{min\_range}, \text{max\_range}]$.

The total number of possible states is: 
$$[\text{Q} (\text{min\_range} + \text{max\_range}) + 1]^{\text{\# network parameters}}$$

\subsection{State Representation and Hashing}
In the A* implementation, the efficacy of the algorithm relies on quickly checking if a particular weight configuration has been visited before, or if a better path to that configuration has been found. This requires a unique, hashable identifier for each state.

The \texttt{QuantizedMLP} utilizes the \texttt{get\_state\_hash()} method, which concatenates all quantized parameters into a single, ordered tuple of integers (by multiplying the quantized weights by the quantization factor $Q$). This tuple serves as the key in the Closed Set (represented by the \texttt{g\_costs} dictionary in the implementation) for tracking the minimum cost found so far to reach that state.

\section{Search Node and Transition Definition}

\subsection{The Search Node}
A single search node $n$ is represented by the \texttt{SearchNode} class and contains:
\begin{enumerate}
    \item The current \texttt{QuantizedMLP} configuration (a unique, quantized set of weights).
    \item The calculated cost-to-come, $g(n)$.
    \item The heuristic estimate to the goal, $h(n)$.
    \item The total estimated cost, $f(n) = g(n) + h(n)$.
\end{enumerate}

\subsection{Defining Transitions (Neighbors)}
The edges of the search graph are defined by the allowed transitions between weight configurations, generated by the \texttt{get\_neighbors} function. A transition represents a minimal, atomic change to the network's weights.

A neighbor state $s$ is generated from the current state $n$ by modifying the value of a single scalar weight by exactly one quantization step ($\Delta w = \pm 1/Q$). The size of the neighborhood, i.e. the total number of neighbors generated for state $n$ is controlled by the parameter $\text{param\_fraction} \in (0, 1)$, which determines the percentage of the total network parameters that are randomly selected to undergo this modification and become a new neighbor state.
$$\text{Transition Rule: } w_{i, s} = w_{i, n} \pm 1/Q$$

This rule ensures that the search space is explored methodically and step-by-step, where the entire training process becomes a controlled traversal of the weight space one minimal adjustment at a time.

\section{A* Cost Function Adaptation for Training}

The successful application of A* to optimization relies entirely on correctly defining the cost components $g(n)$ and $h(n)$ in the context of neural network training.

\subsection{Heuristic Function $h(n)$: Distance to Target Loss}
The goal state is defined as any weight configuration where the network's loss $L(n)$ is less than or equal to a predetermined target loss $L_{\text{target}}$. The heuristic $h(n)$ is defined as the margin by which the current loss exceeds the target loss:

$$h(n) = L(n) - L_{\text{target}}$$

This heuristic serves as the primary driver of the search, prioritizing configurations with lower loss values.

\subsection{Path Cost $g(n)$: Accumulated Iterative Cost}
The cost-to-come, $g(n)$, represents the cost of the path taken from the initial weight configuration to the current node $n$. In this context, it represents the total effort expended in terms of the number of weight updates.

The cost of each transition (edge) is treated as a uniform constant, $\Delta g$, where $\Delta g$ is calculated based on the initial loss and the maximum allowed iterations ($\Delta g = L_{\text{initial}} / I_{\text{max}}$). The total path cost is:

$$g(n) = g(\text{parent}) + \Delta g$$

This definition forces A* to prioritize finding a solution in the fewest number of discrete weight modifications, thereby optimizing for training efficiency in terms of the number of state explorations.

\subsection{Total Evaluation Function $f(n)$}
The combined evaluation function $f(n)$ drives the search priority:

$$f(n) = g(n) + h(n) = \left(g(\text{parent}) + \Delta g\right) + \left(L(n) - L_{\text{target}}\right)$$

By selecting the node with the minimum $f(n)$, the algorithm seeks a balance between minimizing the current error ($h(n)$) and minimizing the total number of weight updates required to reach that state ($g(n)$).

%%%%%%%%%%%%%%%%%%%%%%%%%%%%%%%%%%%%%%%%%%%%%%%%%%%%%%%%%%%%%%%%%%

%\noindent \textbf{Critical Analysis of Admissibility}
\subsection{Admissibility Analysis}

\noindent In standard A* search, $h(n)$ must be admissible ($h(n) \le h^*(n)$, never overestimating the true minimum path cost) to guarantee finding the optimal solution.

\noindent Given the definitions above, the admissibility requirement is not guaranteed and is, in fact, likely violated:
\begin{enumerate}
    \item The heuristic $h(n)$ represents the remaning loss to reach a goal stase, while the unit of the true path cost $h^*(n)$ is a value proportional to the minimum number of parameter modifications ($\Delta g$ steps) required to reach the target loss.
    \item Since the minimum number of steps required to eliminate the remaining loss is unknown a priori and bears no guaranteed mathematical relationship to the loss value itself, $h(n)$ cannot reliably be confirmed as a lower bound for $h^*(n)$.
\end{enumerate}

%\noindent \textbf{Role of $g(n)$ (Non-Optimality):}
\subsection{Role of $g(n)$}

\noindent Since admissibility is ignored, the purpose of $g(n)$ shifts from ensuring global optimality to controlling local exploration. The $g(n)$ term acts both as an iteration counter and as a penalty term, preventing the search from becoming overly greedy and stuck:
\begin{itemize}
    \item It introduces a growing cost proportional to the search effort invested.
    \item This mechanism forces the algorithm to eventually abandon local minima (configurations with low $h(n)$ but high accumulated $g(n)$) and explore alternative configurations on different branches that have been reached with fewer iterations, even if those alternatives currently exhibit a slightly worse $h(n)$ (higher loss).
\end{itemize}
The search thus functions as a Loss-Guided Best-First Search with an Iteration-Based Cost Penalty, trading the A*'s guarantee of optimality for a controlled, non-greedy exploration strategy that helps escape local minima.

%\noindent \textbf{Search Termination}
\subsection{Search Termination}

\noindent Crucially, the search terminates when a goal state is reached, defined by the condition $L(n) \le L_{\text{target}}$, which translates to $h(n) \le 0$. If $L(n) < L_{\text{target}}$, $h(n)$ becomes negative, indicating that the network has surpassed the required goal threshold. Since the search halts immediately upon finding any state where $h(n) \le 0$, the negative heuristic does not represent a practical problem.

%%%%%%%%%%%%%%%%%%%%%%%%%%%%%%%%%%%%%%%%%%%%%%%%%%%%%%%%%%%%%%%%%%%%%%%%%%%%%%%%%%%%%%%%%%%%%%%%%%%%%

\section{The A* Training Loop Implementation}

The core training logic is contained within the \texttt{Trainer} class, which manages the A* search:
\begin{enumerate}
    \item \textbf{Initialization:} The process begins by creating an \texttt{initial\_node} from the starting MLP weights. This node is placed in the priority queue (\texttt{open\_set}).
    \item \textbf{Iteration:} In each iteration, the node with the lowest $f(n)$ is retrieved from the \texttt{open\_set}.
    \item \textbf{Optimality Check:} Before expanding the node, the current path cost $g(n)$ is checked against the minimum cost previously recorded for that state in \texttt{g\_costs} to ensure that A* always processes the optimal path to a given state first.
    \item \textbf{Expansion:} Neighbors are generated by the transition rule ($\pm 1/Q$ change to a weight subset).
    \item \textbf{Cost Update:} For each neighbor, the new $g$ and $h$ values are calculated. If the new path to the neighbor results in a lower cost-to-come ($g_{\text{new}} < g_{\text{old}}$), the neighbor is either updated in the search space or added to the \texttt{open\_set}.
    \item \textbf{Termination:} The search terminates when the best node retrieved from the \texttt{open\_set} satisfies the goal condition ($h(n) \le 0$).
\end{enumerate}

This framework provides a global, informed search mechanism for MLP training, aiming to overcome the local minimum issues inherent in traditional gradient-based methods like Backpropagation.
